\section{Tree of Life Sefirotic Formal Mapping}
\label{sec:tree_of_life_mapping}

We formalize the classical Sefirotic Tree of Life as a directed graph network $\mathcal{G}_{Tree} = (Q, E_{22})$.

\begin{definition}[Sefirotic Node Object]
Each node $Q_i$ ($i \in \{1, \dots, 10\}$) is represented by a structured tuple:
\begin{equation}
Q_i = (M_i, R_i, S_i, V_i, \Phi_i)
\end{equation}
where $M_i$ is station level, $R_i$ is relation set, $S_i$ is scale, $V_i$ is visibility parameter, and $\Phi_i$ is local transition operator.
\end{definition}

\begin{table}[h!]
\centering
\small
\begin{tabular}{c l c c l}
\toprule
\textbf{Node $Q_i$} & \textbf{Traditional Name} & \textbf{Station $M_i$} & \textbf{Scale $S_i$} & \textbf{Category / Role} \\
\midrule
$Q_1$ & Keter & 1 & $1.0$ & Source / Primordial Head \\
$Q_2$ & Chochmah & 2 & $0.9$ & Primary Differentiation \\
$Q_3$ & Binah & 2 & $0.9$ & Structural Receptacle \\
$Q_4$ & Chesed & 3 & $0.7$ & Expansive Flow \\
$Q_5$ & Gevurah & 3 & $0.7$ & Restrictive Bound \\
$Q_6$ & Tiferet & 3 & $0.7$ & Central Harmonizer \\
$Q_7$ & Netzach & 4 & $0.5$ & Active Endurance \\
$Q_8$ & Hod & 4 & $0.5$ & Reflective Echo \\
$Q_9$ & Yesod & 4 & $0.5$ & Foundation Channel \\
$Q_{10}$ & Malchut & 5 & $0.2$ & Physical Manifestation \\
\bottomrule
\end{tabular}
\caption{Mathematical Specification of 10 Sefirotic Nodes}
\end{table}

The 22 connecting channels correspond to edge set $E_{22}$, forming a closed network topology.
